\originalchapter{22.02 作业七:壳层预测}
题源: Spring 2012 Problem Set 7, 2012年4月15日. 壳填充采用课程平均场次序, 预测同实测值区分. 除标记示范外均为编者独立解答.
\begin{exercise}\label{ex:2202-7-1}
原题 1. (a) 给 $^{173}$Yb 基态自旋宇称 (已解). (b) 预测 $^7$Li、$^{15}$N、$^{29}$Si、$^{43}$Ca、$^{87}$Ru、$^{133}$Cs 的基态自旋宇称.
\end{exercise}
\begin{solution}
(a) 官方示范: Yb 有 N=103, 未配对中子在 $2f_{5/2}$, 给 $5/2^-$. (b) 逐核找奇数种类: Li 的 Z=3 在 $1p_{3/2}$, $3/2^-$; N 的 Z=7 在 $1p_{1/2}$, $1/2^-$; Si 的 N=15 在 $2s_{1/2}$, $1/2^+$; Ca 的 N=23 在 $1f_{7/2}$, $7/2^-$; Ru 的 N=43 在课程次序的 $1g_{9/2}$, 简单预测 $9/2^+$; Cs 的 Z=55 在 $1g_{7/2}$, $7/2^+$. $^{87}$Ru 的 NUBASE2020 估计标记为 $1/2^-$, 显示简单球壳预测不等于已测定真值. 不把教材的预测默换成实验表.
\end{solution}
\begin{exercise}\label{ex:2202-7-2}
原题 2. (a) 按课程题称 $^9$Be 的两个激发态 $214$ keV、$1.1$ MeV 给自旋宇称 (已解). (b) $^{81}$Br 基态 $3/2^-$, 三激发态 $5/2^-$ ($275.985$ keV)、$9/2^+$ ($536.20$ keV)、$1/2^-$ ($538.20$ keV), 用壳模型解释四态.
\end{exercise}
\begin{solution}
(a) 原题示范把较低态解释为从 $s_{1/2}$ 壳促升一个核子以留 $s_{1/2}$ 空穴, 给 $1/2^+$; 较高态把价中子移到 $p_{1/2}$, 给 $1/2^-$. 题中 $^9$Be 能量与通行核能级不符, 本稿只保留这套题设模型解释, 不标作实际 $^9$Be 谱. (b) Z=35 的价质子/空穴可在 $2p_{3/2},1f_{5/2},1g_{9/2},2p_{1/2}$ 组态分别给 $3/2^-,5/2^-,9/2^+,1/2^-$. 配对和残余相互作用改变占据能量, 因此仅凭无相互作用填壳顺序不能预测这些精确 keV 值.
\end{solution}
\begin{exercise}\label{ex:2202-7-3}
原题 3. $E_0$ 为无自旋轨道耦合能量, 相互作用 $V_{so}\mathbf L\cdot\mathbf S/\hbar^2$. (a) 两组态能量及差? (b) 证明简并加权重心为 $E_0$ (已解). (c) 填壳预测 $^{209}$Bi 自旋宇称. (d) 哪种作用导致 82 魔数的分裂, Bi 的轨道与 82 壳顶轨道能差多少?
\end{exercise}
\begin{solution}
(a) $E_+=E_0+lV_{so}/2$, $E_-=E_0-(l+1)V_{so}/2$, 差为 $(2l+1)V_{so}/2$. (b) 官方示范: $[(2l+2)E_++2lE_-]/(4l+2)=E_0$. (c) Z=83 超 82 闭壳一个质子, $1h_{9/2}$ 给 $9/2^-$. (d) 强自旋轨道作用使 $1h_{11/2}$ 下移参与 82 闭壳, 同 l=5 的 $h_{9/2},h_{11/2}$ 劈裂大小为 $11|V_{so}|/2$. 原题未给 $V_{so}$ 的实数值, 因而不能虚报一个唯一 MeV 差; 上式是其能量参数表达式, 实际还含其他平均场修正.
\end{solution}
\begin{exercise}\label{ex:2202-7-4}
原题 4. $^{195}$Pt 实际 $1/2^-$ 而简单预测 $13/2^+$, $^{205}$Tl 实际 $1/2^+$ 而预测 $11/2^-$; $^{87}$Sr 基态如预测 $9/2^+$, 有约3小时的 $1/2^-$ 激发态. (a) 哪项核力使前两态更有利, 由 Sr 差别看其依赖? (b) 给三核符合自旋宇称的壳组态.
\end{exercise}
\begin{solution}
配对会让高 j 轨道倾向完整成对, 把未配对核子/空穴留在低 j 近邻轨道. 效果同轨道简并度与配对矩阵元、能隙竞争有关, 不只是所有核恒定的一项. Pt 的 N=117 近126闭壳, 重排后可让 $i_{13/2}$ 核子成对并留下 $p_{1/2}$ 空穴; Tl 的 Z=81 近82闭壳, 可使 $h_{11/2}$ 成对并留下 $s_{1/2}$ 空穴; Sr 的 N=49 近50闭壳, 基态 $g_{9/2}$ 空穴, 激发组态可为 $p_{1/2}$ 空穴. 后者到 $9/2^+$ 需大角动量与宇称改变, 电磁跃迁受抑, 故可能长寿命. 未给组态相互作用参数不能从此模型算出完整寿命.
\end{solution}
\begin{exercise}\label{ex:2202-7-5}
原题 5. $^{40}$Ca 的中子振子模型 $E_N=\hbar\omega N-V'_0$, $V'_0=33$ MeV, 零点能已吸收到井深. (a) 哪些 N 束缚, 首个非束缚 N? (b) 每核子平均结合能? (c) 同质量差所得 $^{40}$Ca 结合能比较.
\end{exercise}
\begin{solution}
采用课程 $\hbar\omega=41A^{-1/3}\simeq11.99$ MeV. $N=0,1,2$ 的单粒子能量约 $-33,-21.01,-9.02$ MeV, N=3 为正 $2.97$ MeV. 每种核子20个占容量 $2,6,12$, 简单单粒子平均 $[2(33)+6(21.01)+12(9.02)]/20\simeq15.02$ MeV. 若把质子势也简化为相同势, 40核子总和约600.7 MeV. AME2020 实际 B 约342.05 MeV, 每核子8.551 MeV, 明显较小. 平均场本征能求和包含相互作用双计数, 还有质子 Coulomb 差别与模型井深误差, 不应把15.02直接称作真实总结合能每核子.
\end{solution}
